\documentclass[11pt, a4paper]{article}
\usepackage[utf8]{inputenc}
\usepackage[T1]{fontenc}
\usepackage[portuguese]{babel}
\usepackage[top=2.5cm, bottom=2.5cm, left=2.5cm, right=2.5cm]{geometry}
\usepackage{enumitem}
\usepackage{xcolor}
\usepackage{tikz}
\usepackage{titlesec}
\usepackage{parskip}
\usepackage{fancyhdr}
\usepackage{lmodern}
\usepackage{microtype}
\usepackage{hyperref}
\usepackage{graphicx}

\hypersetup{hidelinks}

% ── Paleta ───────────────────────────────────────────────────────────────────
\definecolor{primary}{RGB}{22, 82, 168}
\definecolor{primarydark}{RGB}{15, 55, 120}
\definecolor{done}{RGB}{22, 130, 62}
\definecolor{donebg}{RGB}{218, 244, 228}
\definecolor{todo}{RGB}{120, 125, 140}
\definecolor{todobg}{RGB}{246, 246, 249}
\definecolor{rulegray}{RGB}{205, 210, 222}
\definecolor{subseccolor}{RGB}{50, 75, 130}
\definecolor{textdark}{RGB}{25, 30, 48}
\definecolor{stecgray}{RGB}{100, 108, 130}

% ── Checkboxes estilo GitHub ──────────────────────────────────────────────────
\newcommand{\chdone}{%
  \begin{tikzpicture}[baseline=-1.8pt]
    \filldraw[fill=donebg, draw=done, line width=0.75pt, rounded corners=1.5pt]
      (0,0) rectangle (9pt,9pt);
    \draw[done, line width=1.1pt, line cap=round, line join=round]
      (1.8pt,4.2pt) -- (3.8pt,7pt) -- (7.2pt,2pt);
  \end{tikzpicture}%
}

\newcommand{\chtodo}{%
  \begin{tikzpicture}[baseline=-1.8pt]
    \filldraw[fill=todobg, draw=rulegray, line width=0.75pt, rounded corners=1.5pt]
      (0,0) rectangle (9pt,9pt);
  \end{tikzpicture}%
}

% ── Seções (Sprint) ───────────────────────────────────────────────────────────
\titleformat{\section}[block]
  {\large\bfseries\color{primary}}
  {}
  {0pt}
  {\leavevmode{\color{primary}\rule[-0.25em]{3.5pt}{1.15em}}\enspace}
  [{\color{primary}\vspace{2pt}\rule{\linewidth}{0.6pt}}]
\titlespacing{\section}{0pt}{1.8em}{0.8em}

\titleformat{\subsection}[block]{\bfseries\color{subseccolor}\small}{}{0em}{}[%
  {\color{rulegray}\vspace{2pt}\hrule height 0.4pt}\vspace{3pt}]
\titlespacing{\subsection}{0pt}{1em}{0.3em}

% ── Listas ───────────────────────────────────────────────────────────────────
\setlist[itemize]{
  leftmargin = 2em,
  itemsep    = 3pt,
  parsep     = 0pt,
  topsep     = 4pt,
  label      = {}
}

% ── Cabeçalho / rodapé ───────────────────────────────────────────────────────
\pagestyle{fancy}
\fancyhf{}
\lhead{\small\color{stecgray}\textbf{STEC} \color{rulegray}· IACLIN Roadmap}
\rhead{\small\color{stecgray} 2026}
\cfoot{\small\color{rulegray}\thepage}
\renewcommand{\headrulewidth}{0.4pt}
\renewcommand{\headrule}{\color{rulegray}\hrule height 0.4pt}

% ─────────────────────────────────────────────────────────────────────────────
\begin{document}
\color{textdark}
% ─────────────────────────────────────────────────────────────────────────────

% ── Cabeçalho do documento ───────────────────────────────────────────────────
\thispagestyle{empty}

\noindent
% ── Logo: coloque o arquivo da logo na mesma pasta e troque "logo.png" pelo nome correto
\includegraphics[height=46pt]{logo.png}%
\hfill{\small\color{stecgray}\bfseries STEC · \the\year}

\vspace{1.8em}

\noindent{\fontsize{32}{36}\selectfont\bfseries\color{primary} IACLIN}

\vspace{0.3em}
\noindent{\fontsize{15}{18}\selectfont\color{textdark} Roadmap de Desenvolvimento}

\vspace{0.5em}
{\color{rulegray}\hrule height 1pt}
\vspace{0.5em}

\noindent{\small\color{stecgray}
  Documento de acompanhamento das sprints — registra o que já foi entregue
  e o que ainda está em desenvolvimento até a versão MVP.\quad
  \textbf{Prazo de entrega: 21 de julho de 2026.}
}

\vspace{1em}
\noindent
\begin{tikzpicture}[baseline=-1.8pt]
  \filldraw[fill=donebg, draw=done, line width=0.75pt, rounded corners=1.5pt]
    (0,0) rectangle (9pt,9pt);
  \draw[done, line width=1.1pt, line cap=round, line join=round]
    (1.8pt,4.2pt) -- (3.8pt,7pt) -- (7.2pt,2pt);
\end{tikzpicture}
\hspace{3pt}{\small\color{done}\bfseries Concluído}
\hspace{1.8em}
\begin{tikzpicture}[baseline=-1.8pt]
  \filldraw[fill=todobg, draw=rulegray, line width=0.75pt, rounded corners=1.5pt]
    (0,0) rectangle (9pt,9pt);
\end{tikzpicture}
\hspace{3pt}{\small\color{todo}\bfseries Em andamento}

\vspace{1.6em}
{\color{rulegray}\hrule height 0.4pt}
\vspace{1.2em}

% =============================================================================
\section{Sprint 1–2 \ · \ Base do Sistema + Agenda}
\noindent{\small\color{stecgray}\textbf{Responsáveis:} Lucas Pires · Jesus Júnior · Lucas Ferreira}\par\vspace{0.5em}
% =============================================================================

\subsection{Workspace / Multi-clínica}
\begin{itemize}
  \item[\chdone] Separar espaço pessoal do profissional e espaço da clínica
  \item[\chdone] Permitir vínculo do profissional com múltiplas clínicas
  \item[\chdone] Criar opção ``Vincular Clínica'' dentro da plataforma
  \item[\chdone] Garantir isolamento de dados entre clínica e profissional
  \item[\chdone] Remover acesso aos dados da clínica após desvinculação
  \item[\chdone] Sidebar com alternância entre clínicas e espaço pessoal
\end{itemize}

\subsection{Cadastro / Login}
\begin{itemize}
  \item[\chdone] Login e cadastro funcional com prioridade para Google
  \item[\chdone] Obrigatoriedade de CRM/CRO conforme categoria profissional
  \item[\chdone] Complementação cadastral após o primeiro acesso
  \item[\chdone] Separação entre perfil profissional e perfil paciente
  \item[\chdone] Impedir acesso sem validação mínima dos dados profissionais
\end{itemize}

\subsection{Agenda}
\begin{itemize}
  \item[\chdone] Agenda padrão por dias da semana (horário, intervalo e duração)
  \item[\chdone] Bloqueio manual de horários
  \item[\chdone] Status: pendente, confirmado, cancelado e reagendado
  \item[\chdone] Aprovação manual da consulta pelo profissional
  \item[\chdone] Consulta solicitada não aparece como confirmada antes da aprovação
  \item[\chdone] Agenda online pública (ativar/desativar)
  \item[\chdone] Separação entre agenda pessoal, clínica, operadora e agenda online
\end{itemize}

\subsection{Pacientes}
\begin{itemize}
  \item[\chdone] Cadastro, edição e listagem de pacientes
  \item[\chdone] Associação de pacientes às consultas
\end{itemize}

% =============================================================================
\section{Sprint 3 \ · \ Prontuário}
\noindent{\small\color{stecgray}\textbf{Responsável:} Jesus Júnior}\par\vspace{0.5em}
% =============================================================================

\subsection{Ficha e histórico}
\begin{itemize}
  \item[\chdone] Ficha clínica do paciente (criar e editar)
  \item[\chdone] Histórico de consultas associado ao paciente
  \item[\chdone] Prontuário vinculado ao contexto correto: clínica ou espaço pessoal
\end{itemize}

\subsection{Registro de atendimento}
\begin{itemize}
  \item[\chdone] Diagnóstico, conduta/evolução, prescrição, encaminhamento e solicitação de exames
  \item[\chdone] Odontograma básico para fluxo odontológico
\end{itemize}

\subsection{Documentos}
\begin{itemize}
  \item[\chdone] Resumo do atendimento com impressão e exportação em PDF
  \item[\chdone] PDF com nome do paciente, profissional, clínica e logo
\end{itemize}

% =============================================================================
\section{Sprint 4 \ · \ Financeiro}
\noindent{\small\color{stecgray}\textbf{Responsáveis:} Lucas Ferreira · Ryan Victor}\par\vspace{0.5em}
% =============================================================================

\subsection{Controle financeiro básico}
\begin{itemize}
  \item[\chdone] Registrar entradas e saídas
  \item[\chdone] Saldo em caixa visível
  \item[\chdone] Contas a pagar e a receber
  \item[\chdone] Registro manual de pagamento
  \item[\chdone] Separação: financeiro pessoal × financeiro da clínica
  \item[\chdone] Isolamento financeiro entre clínicas
\end{itemize}

\subsection{Funil de vendas / Pipeline}
\begin{itemize}
  \item[\chdone] Pipeline de orçamentos em formato Kanban
  \item[\chdone] Status: aprovado, em negociação e perdido
\end{itemize}

\subsection{Comissionamento de Profissionais}
\begin{itemize}
  \item[\chtodo] Regra de comissão por profissional (percentual ou valor fixo)
  \item[\chtodo] Cálculo de comissão sobre procedimentos realizados no período
  \item[\chtodo] Relatório de comissão por profissional no painel financeiro
  \item[\chtodo] Isolamento de comissão por clínica/workspace
\end{itemize}

% =============================================================================
\section{Sprint 5 \ · \ Integração com Operadoras}
\noindent{\small\color{stecgray}\textbf{Responsáveis:} Lucas Pires · Jesus Júnior}\par\vspace{0.5em}
% =============================================================================

\subsection{Estrutura da operadora}
\begin{itemize}
  \item[\chdone] Painel, autenticação e configuração inicial da operadora
  \item[\chdone] Estrutura separada para operadora médica e odontológica
\end{itemize}

\subsection{Rede credenciada}
\begin{itemize}
  \item[\chdone] Solicitação de vínculo, aceite e recusa
  \item[\chdone] Exibição do status do credenciamento
  \item[\chdone] Listagem de profissionais disponíveis para credenciamento
  \item[\chtodo] Visualizar profissionais por região e especialidade
  \item[\chtodo] Visualizar disponibilidade liberada pelo profissional
\end{itemize}

\subsection{API bidirecional com operadoras — diferencial do produto}
\begin{itemize}
  \item[\chtodo] Conexão bidirecional entre o sistema e o portal da operadora via API
  \item[\chtodo] Sincronização automática da agenda do dentista para a operadora em tempo real
  \item[\chtodo] Atualização imediata dos procedimentos habilitados pelo plano
  \item[\chtodo] Configuração de serviços por operadora (manualmente configurável no MVP)
\end{itemize}

\subsection{Permissões (LGPD)}
\begin{itemize}
  \item[\chdone] Operadora visualiza: agenda liberada, disponibilidade, especialidade e vínculo
  \item[\chdone] Operadora \textbf{não} visualiza: prontuário, exames, diagnóstico e histórico clínico
  \item[\chdone] Operadora médica vê apenas médicos; odontológica vê apenas dentistas
\end{itemize}

% =============================================================================
\section{Sprint 6 \ · \ Credenciamento}
\noindent{\small\color{stecgray}\textbf{Responsáveis:} Lucas Pires · Jesus Júnior}\par\vspace{0.5em}
% =============================================================================

\subsection{Solicitação de vínculo}
\begin{itemize}
  \item[\chtodo] Fluxo de solicitação de vínculo com operadora
  \item[\chtodo] Lista e busca de operadoras (por região e tipo)
  \item[\chtodo] Seleção de serviços/procedimentos na solicitação
  \item[\chtodo] Exibir e registrar status da solicitação com data
\end{itemize}

\subsection{Análise pela operadora}
\begin{itemize}
  \item[\chdone] Operadora visualiza solicitação, dados profissionais e especialidade
  \item[\chdone] Operadora aprova ou recusa com registro da resposta
\end{itemize}

\subsection{Aceite digital e separação por categoria}
\begin{itemize}
  \item[\chtodo] Aceite digital simples com registro de vínculo ativo
  \item[\chtodo] Estrutura preparada para contrato digital futuro
  \item[\chtodo] Separação: credenciamento médico × odontológico
\end{itemize}

% =============================================================================
\section{Sprint 7 \ · \ Marketplace}
\noindent{\small\color{stecgray}\textbf{Responsáveis:} Jesus Júnior · Lucas Ferreira · Ryan Victor}\par\vspace{0.5em}
% =============================================================================

\subsection{Busca e filtros}
\begin{itemize}
  \item[\chdone] Busca por nome, cidade e especialidade
  \item[\chdone] Filtro por localização/geolocalização, distância e tipo de atendimento
  \item[\chdone] Filtro por agenda disponível
  \item[\chdone] Filtro por convênio
\end{itemize}

\subsection{Perfil do profissional}
\begin{itemize}
  \item[\chdone] Especialidades, localização e disponibilidade conforme agenda
  \item[\chdone] Exibir apenas dados públicos permitidos
  \item[\chtodo] Foto do profissional
  \item[\chtodo] Convênios aceitos exibidos no perfil
\end{itemize}

\subsection{Agenda online}
\begin{itemize}
  \item[\chdone] Exibir apenas profissionais com agenda online ativada
  \item[\chdone] Solicitação de agendamento com sinalização ``aguardando confirmação''
  \item[\chdone] Respeitar agenda pessoal, clínica e operadora
\end{itemize}

% =============================================================================
\section{Sprint 8 \ · \ Comunicação e Automação}
\noindent{\small\color{stecgray}\textbf{Responsáveis:} Lucas Pires · Jesus Júnior · Endria Carem}\par\vspace{0.5em}
% =============================================================================

\subsection{WhatsApp}
\begin{itemize}
  \item[\chdone] Fluxo de conexão do WhatsApp com a IA Secretária
  \item[\chtodo] Notificações de consulta aprovada, recusada e reagendada
  \item[\chtodo] Lembrete manual de consulta
  \item[\chtodo] Mensagens padrão de confirmação, reagendamento e cancelamento
\end{itemize}

\subsection{IA Secretária}
\begin{itemize}
  \item[\chdone] Prompt base estruturado
  \item[\chdone] Consultar disponibilidade real da agenda
  \item[\chdone] Não oferecer horário indisponível
  \item[\chdone] Encaminhamento para atendimento humano (handoff)
  \item[\chdone] Responder mensagens de pacientes automaticamente
  \item[\chtodo] Auxiliar no agendamento perguntando se é particular ou plano
  \item[\chtodo] Informar serviços, convênios aceitos e localização da clínica
  \item[\chtodo] Informar status da consulta (inclusive quando aguarda aprovação)
  \item[\chtodo] Respeitar agenda de plano/operadora
\end{itemize}

\subsection{IA Gestor}
\begin{itemize}
  \item[\chdone] Chat interno disponível para clínica e profissional autônomo
  \item[\chdone] Responder sobre consultas do dia, próximas, pendentes e canceladas
  \item[\chdone] Responder sobre horários disponíveis e volume de atendimentos
  \item[\chdone] Base preparada para análise financeira futura
\end{itemize}

\subsection{Automações}
\begin{itemize}
  \item[\chdone] Estrutura de automações (criar, editar e ativar regras de disparo)
  \item[\chtodo] Fluxo de lembrete de consulta (24h antes)
  \item[\chtodo] Fluxo de confirmação automática
  \item[\chtodo] Fluxo de reagendamento
  \item[\chtodo] Fluxo de retorno do paciente
\end{itemize}

\subsection{Campanhas de WhatsApp}
\begin{itemize}
  \item[\chtodo] Envio em massa de mensagens via WhatsApp
  \item[\chtodo] Template de mensagem personalizável por campanha
  \item[\chtodo] Filtrar base de pacientes para envio (todos, por status, por período)
  \item[\chtodo] Relatório básico de envio (enviados / falhos)
\end{itemize}

\subsection{Régua de Relacionamento}
\begin{itemize}
  \item[\chtodo] Mensagem automática de aniversário para o paciente
  \item[\chtodo] Lembrete de retorno preventivo (ex: ``Faz 6 meses desde sua limpeza'')
  \item[\chtodo] Pesquisa de satisfação (NPS) automática pós-consulta
  \item[\chtodo] Gatilhos por data de aniversário, última consulta e consulta realizada
\end{itemize}

% =============================================================================
\section{Sprint 9 \ · \ Finalização e Deploy}
\noindent{\small\color{stecgray}\textbf{Responsáveis:} Lucas Pires · Jesus Júnior · Lucas Ferreira · Ryan Victor}\par\vspace{0.5em}
% =============================================================================

\subsection{Cobrança de Assinaturas (Stripe)}
\begin{itemize}
  \item[\chtodo] Planos de assinatura no Stripe (Clínica / Profissional Autônomo)
  \item[\chtodo] Checkout de assinatura integrado
  \item[\chtodo] Status da assinatura ativa visível no painel
  \item[\chtodo] Bloqueio de acesso em caso de assinatura vencida ou cancelada
  \item[\chtodo] Upgrade, downgrade e cancelamento de plano
  \item[\chtodo] Webhook do Stripe para atualizar status em tempo real
\end{itemize}

\subsection{Validação final e QA}
\begin{itemize}
  \item[\chtodo] Validar agenda, marketplace, operadoras e fluxos de IA
  \item[\chtodo] Garantir que dados clínicos não fiquem expostos
  \item[\chtodo] Responsividade mobile e ajustes visuais finais
  \item[\chtodo] Revisar e corrigir bugs críticos restantes
\end{itemize}

\subsection{Deploy}
\begin{itemize}
  \item[\chtodo] Deploy do sistema em ambiente online
  \item[\chtodo] Validar ausência de erros críticos em produção
\end{itemize}

% =============================================================================
\section{Sprint 10 \ · \ Especialidades}
\noindent{\small\color{stecgray}\textbf{Responsáveis:} Lucas Pires · Jesus Júnior · Lucas Ferreira}\par\vspace{0.5em}
% =============================================================================

\subsection{Estrutura padrão reutilizável}
\begin{itemize}
  \item[\chtodo] Validar fluxo padrão reutilizável entre especialidades
  \item[\chtodo] Identificar campos realmente específicos de cada área
  \item[\chtodo] Odontologia: validar odontograma básico e fluxo específico
\end{itemize}

\subsection{Especialidades médicas}
\begin{itemize}
  \item[\chtodo] Validar necessidades específicas: pediatria, psicologia, ortopedia, dermatologia e cardiologia
  \item[\chtodo] Identificar quais especialidades precisam de telas próprias no MVP
\end{itemize}

% ── Rodapé do documento ───────────────────────────────────────────────────────
\vspace{2em}
{\color{rulegray}\hrule height 0.4pt}
\vspace{0.6em}
\noindent{\small\color{stecgray}
  \textbf{STEC} · Documento de uso interno e apresentação ao cliente · \the\year
}

\end{document}
